% Apunts de teoria en PDF · 4t ESO · Unitat 11: Probabilitat
% Mateix estil que els apunts de Batxillerat (franja i caixes de teoria en blau, exemples en roig), però més visuals,
% curts i amb dibuixos, per a l'alumnat de 3r ESO. Es compila amb pdflatex dues vegades (per als enllaços de l'índex),
% amb logo-fentmates.pdf (_eines/logo/pdf/horitzontal-clar.pdf) al costat.
\documentclass[11pt,a4paper]{article}
\usepackage[utf8]{inputenc}
\usepackage[T1]{fontenc}
\usepackage[catalan]{babel}
\usepackage[sfdefault,book,lining]{FiraSans}
\usepackage[italic,noplusnominus,noequal]{mathastext}
\linespread{1.18}
\usepackage[a4paper,top=2.4cm,bottom=2.2cm,left=1.6cm,right=1.6cm,headheight=1.1cm,headsep=0.5cm,footskip=1cm]{geometry}
\usepackage{amsmath,amssymb,cancel}
\usepackage{icomma}
\usepackage{graphicx}
\usepackage[table]{xcolor}
\usepackage{tikz}\usetikzlibrary{calc,positioning,arrows.meta,decorations.pathreplacing,shapes.geometric}
\usepackage{fontawesome5}
\usepackage{ccicons}
\usepackage{fancyhdr}
\usepackage{enumitem}
\usepackage{array}
\usepackage{needspace}
\usepackage[most]{tcolorbox}
\usepackage{titlesec}
\usepackage[hidelinks]{hyperref}
\hypersetup{pdftitle={Apunts de teoria. Unitat 11: Probabilitat. 4t ESO},pdfsubject={Matemàtiques A 4t ESO},pdfkeywords={experiments aleatoris, espai mostral, esdeveniments, unió, intersecció, contrari, freqüència relativa, llei dels grans nombres, llei de Laplace, experiències compostes, regla del producte, probabilitat condicionada, probabilitat total, taules de contingència}}

% colors de la web (bloc d'Àlgebra: blau)
\definecolor{fmnavy}{HTML}{1B2748}
\definecolor{fmblue}{HTML}{3D6AA0}
\definecolor{fmbluel}{HTML}{4E81BD}
\definecolor{fmyellow}{HTML}{FFC94D}
\definecolor{fmgold}{HTML}{E9A919}
\definecolor{fmaccent}{HTML}{FF7A50}
\definecolor{fmdef}{HTML}{1D4F73}\definecolor{fmdefbg}{HTML}{EEF4F9}
\definecolor{fmex}{HTML}{A3203A}\definecolor{fmexbg}{HTML}{FBF0F2}
\definecolor{fmsf}{HTML}{C2410C}
\definecolor{fmok}{HTML}{15803D}
\definecolor{fmmuted}{HTML}{5D6685}
\definecolor{fmsurface}{HTML}{EEF1F8}
\definecolor{fmA}{HTML}{2F6FD6}\definecolor{fmB}{HTML}{EA6A10}\definecolor{fmR}{HTML}{E11D48}\definecolor{fmK}{HTML}{16A34A}\definecolor{fmV}{HTML}{7C3AED}

% lletres
\newcommand{\comic}{\fontfamily{ComicNeue-TLF}\selectfont}
\newcommand{\heavy}{\fontseries{eb}\selectfont}
\newcommand{\semi}{\fontseries{sb}\selectfont}
\DeclareUnicodeCharacter{2212}{-}

\setlength{\parindent}{0pt}
\setlength{\parskip}{7pt}
\setlength{\lineskip}{4pt}\setlength{\lineskiplimit}{4pt}
\setlength{\emergencystretch}{3em}
\setlist{itemsep=3pt,topsep=3pt,leftmargin=1.4em}
\setcounter{secnumdepth}{1}
\renewcommand{\thesection}{11.\arabic{section}}
\titleformat{\section}[hang]{\color{fmnavy}\fontsize{23}{27}\heavy\selectfont}{\tcbox[on line,colback=fmblue,colframe=fmblue,arc=3mm,boxsep=2pt,left=6pt,right=6pt,top=3pt,bottom=3pt]{\color{white}\thesection}}{0.45em}{}[\vspace{5pt}{\color{fmblue}\titlerule[2pt]}]
\titlespacing*{\section}{0pt}{4pt}{6pt}
\titleformat{\subsection}{\color{fmblue}\Large\semi}{}{0pt}{\faCaretRight\ }
\titlespacing*{\subsection}{0pt}{8pt}{2pt}

% frase de cada apartat, «a mà»
\newcommand{\frase}[1]{\par\vspace{-2pt}{\comic\Large\color{fmbluel}\faCommentDots[regular]\ \ #1\par}\vspace{2pt}}
\newcommand{\seccio}[3]{\Needspace{16\baselineskip}\section{#2}\hypertarget{#1}{}\frase{#3}}
% notes a mà dins dels exemples
\newcommand{\nota}[1]{{\comic\small\color{fmex!85!black}#1}}
\newcommand{\notab}[1]{{\comic\small\color{fmblue}#1}}
% bé / malament
\newcommand{\be}[1]{{\color{fmok}\faCheckCircle}\ \ensuremath{#1}}
\newcommand{\mal}[1]{{\color{fmR}\faTimesCircle}\ {\color{fmR}\ensuremath{#1}}}
% pas numerat
\newcommand{\pas}[1]{\tikz[baseline=(p.base)]\node[circle,fill=fmblue,text=white,inner sep=1.2pt,minimum size=1.25em,font=\small\heavy](p){#1};}
% etiqueta de color
\newcommand{\chip}[2]{\tcbox[on line,colback=#1!14,colframe=#1!60,arc=2mm,boxsep=1pt,left=5pt,right=5pt,top=1pt,bottom=1pt]{#2}}
\newcommand{\fmkey}[1]{\colorbox{fmyellow!55}{\strut #1}}
\newcommand{\cA}[1]{{\color{fmA}#1}}\newcommand{\cB}[1]{{\color{fmB}#1}}\newcommand{\cK}[1]{{\color{fmK}#1}}\newcommand{\cR}[1]{{\color{fmR}#1}}\newcommand{\cV}[1]{{\color{fmV}#1}}

% caixes
\tcbset{fmbase/.style={enhanced,arc=3mm,boxrule=1.2pt,left=4mm,right=4mm,top=2mm,bottom=2.5mm,before skip=8pt,after skip=8pt,fonttitle=\heavy\large,before upper=\raggedright}}
\newtcolorbox{idea}[1][Idea clau]{fmbase,colback=fmdefbg,colframe=fmdef,coltitle=white,title={\faKey\ \ #1}}
\newtcolorbox{exemple}[1][Exemple]{fmbase,colback=fmexbg,colframe=fmex,colbacktitle=fmexbg,coltitle=fmex,title={\faPencil*\ \ #1},titlerule=0.5pt,titlerule style={fmex!40,dashed}}
\newtcolorbox{compte}[1][Compte!]{fmbase,colback=fmsf!6,colframe=fmsf,coltitle=white,fonttitle=\comic\bfseries\Large,title={\faExclamationTriangle\ \ #1}}
\newtcolorbox{truc}[1][Truc]{fmbase,colback=fmyellow!20,colframe=fmgold,colbacktitle=fmyellow,coltitle=fmnavy,fonttitle=\comic\bfseries\Large,title={\faLightbulb[regular]\ \ #1}}

% dibuix de la targeta de la unitat, en blanc
\newcommand{\fmcard}{\begin{tikzpicture}[x=0.062cm,y=-0.062cm,white,line width=3pt,line cap=round]
  \begin{scope}[transparency group,opacity=0.15]
  \node[anchor=base west,inner sep=0pt,font=\fontsize{28}{28}\selectfont] at (14,44) {$P(A)\quad\frac{3}{6}$};
  \draw (140,70) -- (180,45) (140,70) -- (180,95);
  \draw (180,45) -- (225,32) (180,45) -- (225,58) (180,95) -- (225,82) (180,95) -- (225,108);
  \foreach \p in {(140,70),(180,45),(180,95),(225,32),(225,58),(225,82),(225,108)} \fill \p circle (5);
  \end{scope}
\end{tikzpicture}}
\newcommand{\fmbrand}{\raisebox{-0.25\height}{\includegraphics[height=0.5cm]{logo-fentmates.pdf}}}

% regla de Ruffini: #1 = a, #2 = coeficients, #3 = fila del mig (sense el primer), #4 = fila de baix, #5 = nombre de coeficients
\newcommand{\ruffini}[5]{\begin{tikzpicture}[x=1.25cm,y=0.85cm,every node/.style={font=\large}]
  \foreach \v [count=\i] in {#2} \node (t\i) at (\i,0) {$\v$};
  \foreach \v [count=\i] in {#3} {\pgfmathtruncatemacro\j{\i+1}\node (m\j) at (\j,-1) {$\v$};}
  \foreach \v [count=\i] in {#4} \node (b\i) at (\i,-2.1) {$\v$};
  \node[text=fmB] at (0,-1) {$#1$};
  \draw[line width=1.1pt,fmnavy] (0.45,0.5)--(0.45,-1.5) (0.45,-1.5)--(#5+0.5,-1.5);
  \node[draw=fmR,line width=1.3pt,rounded corners=2mm,minimum size=0.75cm] at (#5,-2.1) {};
  \draw[-{Latex[length=2.2mm]},fmK,line width=1.1pt] (t1.south) -- (b1.north);
  \pgfmathtruncatemacro\last{#5-1}
  \foreach \i in {1,...,\last} {\pgfmathtruncatemacro\j{\i+1}
    \draw[-{Latex[length=2.2mm]},fmB,line width=1.1pt,shorten >=1pt] (b\i.north east) -- (m\j.south west);
    \draw[-{Latex[length=2.2mm]},fmA!70,line width=1pt,dashed,shorten >=2pt] ([xshift=8pt]t\j.south) -- ([xshift=8pt]b\j.north);}
\end{tikzpicture}}

% capçalera i peu
\pagestyle{fancy}\fancyhf{}
\fancyhead[L]{\small\fmbrand}
\fancyhead[R]{\small\color{fmmuted}4t ESO · Unitat 11: Probabilitat}
\fancyfoot[L]{\scriptsize\color{fmmuted}\href{https://creativecommons.org/licenses/by-nc-sa/4.0/deed.ca}{\ccbyncsa}\ \ FentMates · fentmates.org · CC BY-NC-SA 4.0}
\fancyfoot[R]{\small\color{fmmuted}\thepage}
\renewcommand{\headrulewidth}{0.4pt}\renewcommand{\headrule}{\color{fmmuted!40}\hrule height \headrulewidth}
\fancypagestyle{primera}{\fancyhf{}\renewcommand{\headrulewidth}{0pt}
  \fancyfoot[L]{\scriptsize\color{fmmuted}\href{https://creativecommons.org/licenses/by-nc-sa/4.0/deed.ca}{\ccbyncsa}\ \ FentMates · fentmates.org · CC BY-NC-SA 4.0}
  \fancyfoot[R]{\small\color{fmmuted}\thepage}}

% targeta de l'índex
\newcommand{\itile}[4]{\begin{tcolorbox}[enhanced,colback=white,colframe=fmblue!35,arc=3mm,boxrule=1pt,boxsep=0pt,left=6pt,right=6pt,top=5pt,bottom=5pt,before skip=0pt,after skip=7pt]
  \hyperlink{#1}{\begin{minipage}[c]{0.9cm}\centering{\color{fmblue}\LARGE #3}\end{minipage}\hspace{6pt}%
  \begin{minipage}[c]{0.78\linewidth}{\color{fmblue}\heavy #2}\\[-2pt]{\color{fmnavy}\semi\large #4}\end{minipage}}\end{tcolorbox}}

% taula amb vores arrodonides: el contingut es retalla amb la forma (sense @{} a les vores ni entre columnes)
\newsavebox\tabclipbox
\newcommand{\taulaclip}[3][3mm]{\sbox\tabclipbox{#3}\begin{tikzpicture}[baseline=(tcn.base)]
  \begin{scope}\clip[rounded corners=#1] (0,-\dp\tabclipbox) rectangle (\wd\tabclipbox,\ht\tabclipbox);
  \node[inner sep=0pt,outer sep=0pt,anchor=base west] (tcn) at (0,0) {\usebox\tabclipbox};\end{scope}
  \draw[#2,rounded corners=#1] (0,-\dp\tabclipbox) rectangle (\wd\tabclipbox,\ht\tabclipbox);
\end{tikzpicture}}
\newcommand{\taularod}[3][fmblue]{\taulaclip{#1,line width=1.2pt}{\renewcommand{\arraystretch}{1.5}\begin{tabular}{#2}#3\end{tabular}}}
% balança
\newcommand{\balanca}[3]{\begin{tikzpicture}[scale=#3,transform shape,every node/.style={anchor=center}]
  \fill[fmnavy] (-0.45,-1.9) rectangle (0.45,-1.75);
  \fill[fmnavy!80] (-0.12,-1.75) rectangle (0.12,0);
  \fill[fmgold] (0,0.05) circle (0.13);
  \draw[fmnavy,line width=2.4pt,line cap=round] (-2,0) -- (2,0);
  \foreach \s in {-1,1} {
    \draw[fmmuted,line width=0.8pt] (\s*2,0) -- (\s*2-0.75,-1.05) (\s*2,0) -- (\s*2+0.75,-1.05);
    \fill[fmyellow!70,draw=fmgold,line width=1pt] (\s*2-0.95,-1.05) arc[start angle=180,end angle=360,x radius=0.95,y radius=0.32] -- cycle;}
  \node[font=\Large,fill=fmA!15,draw=fmA,rounded corners=1.5mm,inner sep=3pt] at (-2,-0.68) {$#1$};
  \node[font=\Large,fill=fmB!15,draw=fmB,rounded corners=1.5mm,inner sep=3pt] at (2,-0.68) {$#2$};
\end{tikzpicture}}
% eixos amb quadrícula (dins d'un tikzpicture): xmin, xmax, ymin, ymax
\newcommand{\eixos}[4]{\draw[fmmuted!22,line width=0.4pt] (#1,#3) grid (#2,#4);
  \draw[-{Latex[length=2mm]},fmmuted,line width=0.8pt] (#1-0.2,0)--(#2+0.4,0) node[right,font=\small]{$x$};
  \draw[-{Latex[length=2mm]},fmmuted,line width=0.8pt] (0,#3-0.2)--(0,#4+0.4) node[above,font=\small]{$y$};}
% taula de traduccions (diccionari) amb les vores arrodonides
\newcommand{\dic}[2]{#1 & {\large$#2$}\\}
\newcommand{\dicratio}{0.62}
\newcommand{\dichead}{\begin{tabular}{@{}>{\raggedright}p{\dimexpr\dicratio\linewidth-12pt}>{\centering\arraybackslash}p{\dimexpr\linewidth-\dicratio\linewidth-12pt}@{}}\faComment[regular]\ Ho diem així\dots & \dots\ així\end{tabular}}
\newenvironment{dicbox}{\begin{tcolorbox}[enhanced,arc=3mm,boxrule=1.2pt,colframe=fmblue,colback=white,coltitle=white,fonttitle=\heavy,
  title={\dichead},left=0pt,right=0pt,top=0pt,bottom=1.5mm,boxsep=0pt,lefttitle=6pt,righttitle=6pt,before skip=0pt,after skip=0pt]\rowcolors{1}{white}{fmsurface}%
  \begin{tabular}{@{\hspace{6pt}}>{\raggedright}p{\dimexpr\dicratio\linewidth-12pt}>{\centering\arraybackslash}p{\dimexpr\linewidth-\dicratio\linewidth-12pt}@{\hspace{6pt}}}}
  {\end{tabular}\end{tcolorbox}}
% resum
\newcommand{\rtile}[3]{\begin{tcolorbox}[enhanced,colback=white,colframe=fmblue!45,arc=3mm,boxrule=1pt,left=3mm,right=3mm,top=2.5mm,bottom=2.5mm,fontupper=\large,
  title={\tcbox[on line,colback=fmblue,colframe=fmblue,arc=2mm,boxsep=0pt,left=4pt,right=4pt,top=1pt,bottom=1pt]{\color{white}\small\heavy #1}\ \ #2},coltitle=fmnavy,colbacktitle=fmdefbg,fonttitle=\semi,toptitle=0.5mm,bottomtitle=0.5mm]#3\end{tcolorbox}}
% portada: #1 subtítol a mà
\newcommand{\portada}[1]{\thispagestyle{primera}
\begin{tikzpicture}[remember picture,overlay]
  \fill[fmblue] (current page.north west) rectangle ([yshift=-9.4cm]current page.north east);
  \begin{scope}
    \clip (current page.north west) rectangle ([yshift=-9.4cm]current page.north east);
    \node[anchor=north west,inner sep=0pt] at ([xshift=6.6cm,yshift=-1.3cm]current page.north west) {\fmcard};
  \end{scope}
  \node[anchor=north west,inner sep=0pt] at ([xshift=1.6cm,yshift=-1.1cm]current page.north west) {\tcbox[colback=white,colframe=white,arc=2.5mm,boxsep=0pt,left=7pt,right=7pt,top=5pt,bottom=5pt]{\includegraphics[height=0.75cm]{logo-fentmates.pdf}}};
\end{tikzpicture}
\vspace*{0.55cm}
\tcbox[colback=white!15!fmblue,colframe=white!15!fmblue,arc=2.5mm,boxsep=0pt,left=6pt,right=6pt,top=3pt,bottom=3pt]{\color{white}\small\bfseries Apunts de teoria · 4t ESO · Matemàtiques A}

\vspace{2pt}
{\raggedright\color{white}\fontsize{30}{36}\heavy\selectfont{\color{fmyellow}Unitat 11}\\[2pt]Probabilitat\par}
\vspace{8pt}
{\color{white}\comic\Large #1\par}
\vspace*{\dimexpr 7.4cm-\pagetotal\relax}
}
\newcommand{\llegenda}{\par\vspace{16pt}
{\color{fmnavy}\heavy\Large\faPalette\ \ Com es llegeixen aquests apunts}\par\smallskip
\begin{tcbraster}[raster columns=4,raster equal height,raster column skip=8pt,raster before skip=12pt]
\tcbox[colback=fmdefbg,colframe=fmdef,arc=3mm,boxrule=1.2pt,halign=center,width=\linewidth]{{\color{fmdef}\Large\faKey}\\[2pt]{\color{fmdef}\heavy Idea clau}\\[-1pt]{\small el que cal saber}}
\tcbox[colback=fmexbg,colframe=fmex,arc=3mm,boxrule=1.2pt,halign=center,width=\linewidth]{{\color{fmex}\Large\faPencil*}\\[2pt]{\color{fmex}\heavy Exemple}\\[-1pt]{\small pas a pas}}
\tcbox[colback=fmsf!6,colframe=fmsf,arc=3mm,boxrule=1.2pt,halign=center,width=\linewidth]{{\color{fmsf}\Large\faExclamationTriangle}\\[2pt]{\color{fmsf}\heavy Compte!}\\[-1pt]{\small errors típics}}
\tcbox[colback=fmyellow!20,colframe=fmgold,arc=3mm,boxrule=1.2pt,halign=center,width=\linewidth]{{\color{fmgold!80!black}\Large\faLightbulb[regular]}\\[2pt]{\color{fmnavy}\heavy Truc}\\[-1pt]{\small per a recordar-ho}}
\end{tcbraster}}
\newcommand{\resumtitol}{\clearpage
\hypertarget{resum}{}
{\color{fmnavy}\fontsize{23}{27}\heavy\selectfont\faClipboardList\ \ En resum}\par\vspace{5pt}{\color{fmblue}\titlerule[2pt]}
\frase{Tota la unitat en una pàgina.}}
% els 5 passos per a resoldre problemes: #2 = text del pas 2
\newcommand{\passosproblema}[2]{\begin{center}
\begin{tikzpicture}[s/.style={draw=fmblue,fill=fmdefbg,rounded corners=3mm,line width=1.2pt,minimum width=2.75cm,minimum height=2.45cm,align=center,font=\small},
  ic/.style={font=\LARGE,text=fmblue}]
  \foreach \i/\ic/\t in {1/\faBookOpen/{\textbf{Llegim}\\i entenem},2/\faQuestionCircle/{#1},3/\faEquals/{#2},4/\faCogs/{\textbf{Resolem}},5/\faCheckDouble/{\textbf{Comprovem}\\i responem}} {
    \node[s] (s\i) at (3.3*\i,0) {};
    \node[ic] at ([yshift=-0.5cm]s\i.north) {\ic};
    \node[align=center,font=\small,anchor=north] at ([yshift=-0.95cm]s\i.north) {\t};
    \node[circle,fill=fmblue,text=white,font=\small\heavy,inner sep=1.5pt,minimum size=1.3em] at (s\i.north west) {\i};}
  \foreach \i [evaluate=\i as \j using int(\i+1)] in {1,...,4} \draw[-{Latex[length=3mm]},fmmuted,line width=1.5pt] (s\i) -- (s\j);
\end{tikzpicture}
\end{center}}

\newcommand{\lo}{\mathopen{]}}\newcommand{\ro}{\mathclose{[}}
% punts ple i buit, nombres dels eixos i taula de valors
\newcommand{\pple}[2][fmblue]{\fill[#1] #2 circle (3.6pt);}
\newcommand{\nums}[4]{\foreach \k in {#1,...,#2} {\ifnum\k=0\else\node[below,font=\scriptsize,text=fmmuted] at (\k,0) {$\k$};\fi}
  \foreach \k in {#3,...,#4} {\ifnum\k=0\else\node[left,font=\scriptsize,text=fmmuted] at (0,\k) {$\k$};\fi}}
\newcommand{\tv}[3][fmblue]{\taulaclip{#1,line width=1.2pt}{\renewcommand{\arraystretch}{1.45}%
  \begin{tabular}{>{\columncolor{#1}\color{white}\heavy}c*{12}{w{c}{0.8cm}}}#2\\\rowcolor{#1!8}\cellcolor{#1}#3\end{tabular}}}
% etiqueta d'una recta o corba
\newcommand{\etq}[3]{\node[fill=white,inner sep=1.5pt,text=#1,font=\small] at #2 {#3};}


\def\intvdot#1#2#3{\ifx#2o\filldraw[fill=white,draw=#3,line width=1.3pt] (#1,0) circle (3.3pt);\fi\ifx#2c\fill[#3] (#1,0) circle (3.5pt);\fi}
\newcommand{\intv}[5]{\begin{tikzpicture}[x=0.4cm,baseline=-3pt]
  \draw[fmmuted,line width=0.8pt] (-6,0)--(7,0);
  \foreach \k in {-5,...,6} \draw[fmmuted,line width=0.6pt] (\k,-0.08)--(\k,0.08);
  \ifx#3a\def\xa{-6}\else\def\xa{#1}\fi
  \ifx#4a\def\xb{7}\else\def\xb{#2}\fi
  \ifx#4a\draw[#5,line width=2.6pt,-{Latex[length=2.6mm]}] (\xa,0)--(\xb,0);\else\ifx#3a\draw[#5,line width=2.6pt,{Latex[length=2.6mm]}-] (\xa,0)--(\xb,0);\else\draw[#5,line width=2.6pt] (\xa,0)--(\xb,0);\fi\fi
  \intvdot{#1}{#3}{#5}\intvdot{#2}{#4}{#5}
  \ifx#3a\else\node[below=2pt,font=\scriptsize] at (#1,0) {$#1$};\fi
  \ifx#4a\else\node[below=2pt,font=\scriptsize] at (#2,0) {$#2$};\fi
\end{tikzpicture}}
\tikzset{br/.style={fmmuted,line width=1pt},nd/.style={circle,fill=fmblue,text=white,font=\scriptsize\heavy,inner sep=1.5pt,minimum size=1.35em},
  pl/.style={font=\scriptsize,fill=white,inner sep=1pt}}
\newcommand{\cond}{\,|\,}
\newcommand{\venn}[2]{\begin{tikzpicture}[scale=0.5]
  \def\ca{(-0.55,0) circle (1)}\def\cb{(0.55,0) circle (1)}
  #1
  \draw[fmnavy,line width=0.8pt] (-1.9,-1.3) rectangle (1.9,1.3); \node[font=\scriptsize,text=fmnavy] at (1.6,1.0) {$E$};
  #2
\end{tikzpicture}}
\newcommand{\bola}[1]{\tikz[baseline=-0.6ex]\shade[ball color=#1] (0,0) circle (0.17cm);}
\tikzset{arbre/.style={grow=right,level distance=2.1cm,edge from parent/.style={draw,fmnavy,line width=0.9pt},every node/.style={font=\small}},
  ram/.style={midway,fill=white,inner sep=1pt,font=\scriptsize,text=fmB}}


\begin{document}
\portada{No sabem què eixirà, però sí com de probable és.}

{\color{fmnavy}\heavy\Large\faMapSigns\ \ Què aprendràs?}\par\smallskip
\begin{minipage}[t]{0.485\linewidth}
\itile{sec1}{11.1}{\faDiceThree}{Experiments}
\itile{sec2}{11.2}{\faIcon{object-group}}{Esdeveniments}
\itile{sec3}{11.3}{\faChartLine}{Probabilitat d'un esdeveniment}
\itile{sec4}{11.4}{\faBalanceScale}{Càlcul de probabilitats}
\end{minipage}\hfill
\begin{minipage}[t]{0.485\linewidth}
\itile{sec5}{11.5}{\faIcon{sitemap}}{Experiències compostes}
\itile{sec6}{11.6}{\faShoppingBag}{Dependents. Condicionada}
\itile{sec7}{11.7}{\faTable}{Taules de contingència}
\itile{resum}{\phantom{11.7}}{\faClipboardList}{En resum}
\end{minipage}
\llegenda

% ============================================================ 11.1
\clearpage
\seccio{sec1}{Experiments}{Espai mostral i esdeveniments.}

\vfill
\begin{tcbraster}[raster columns=2,raster equal height,raster column skip=10pt]
\begin{tcolorbox}[fmbase,colback=fmB!8,colframe=fmB,coltitle=white,title={\faCogs\ \ Determinista}]
Sabem el resultat abans de fer-lo.\\ \small Ex.: el temps que tarda a caure una pilota.
\end{tcolorbox}
\begin{tcolorbox}[fmbase,colback=fmK!8,colframe=fmK,coltitle=white,title={\faDiceThree\ \ Aleatori}]
No el podem predir, però coneixem els possibles.\\ \small Ex.: llançar un dau.
\end{tcolorbox}
\end{tcbraster}

\vfill
\begin{idea}[Espai mostral i esdeveniments]
\textbf{Espai mostral} $E$: tots els resultats possibles. \ \textbf{Esdeveniment}: qualsevol subconjunt de $E$ ($A$, $B$\dots). Passa si el resultat és un dels seus elements.
\end{idea}

\vfill
\begin{tcbraster}[raster columns=4,raster equal height,raster column skip=6pt]
\begin{tcolorbox}[fmbase,colback=fmdefbg,colframe=fmdef,coltitle=white,title={Elemental}]
\small Un sol resultat: $\{5\}$
\end{tcolorbox}
\begin{tcolorbox}[fmbase,colback=fmdefbg,colframe=fmdef,coltitle=white,title={Compost}]
\small Dos o més: $\{2,4,6\}$
\end{tcolorbox}
\begin{tcolorbox}[fmbase,colback=fmdefbg,colframe=fmdef,coltitle=white,title={Segur}]
\small $E$: sempre passa
\end{tcolorbox}
\begin{tcolorbox}[fmbase,colback=fmdefbg,colframe=fmdef,coltitle=white,title={Impossible}]
\small $\varnothing$: no passa mai
\end{tcolorbox}
\end{tcbraster}

\vfill
\begin{exemple}[Exemple: \ llançar un dau]
$E=\{1,2,3,4,5,6\}$ \qquad «un 5» $=\{5\}$ elemental \qquad «parell» $=\{2,4,6\}$ compost\\[4pt]
«menor que 7» $=E$ segur \qquad «un 8» $=\varnothing$ impossible
\end{exemple}

\vfill
\begin{exemple}[Exemple: \ urna amb boles de l'1 al 12]
$A=$ «parell» $=\{2,4,6,8,10,12\}$ \qquad $C=$ «múltiple de 3» $=\{3,6,9,12\}$ \qquad $D=$ «múltiple de 6» $=\{6,12\}$
\end{exemple}

\vfill
\begin{exemple}[Exemples]
Tres monedes: \ $2\cdot2\cdot2=8$ resultats: CCC, CCX, CXC, CXX, XCC, XCX, XXC, XXX\\[4pt]
Dos daus: \ $6\cdot6=36$ parelles \qquad Família de 3 fills, nombre de xics: \ $E=\{0,1,2,3\}$\\[4pt]
Pal d'una carta: \ $E=\{\text{oros, copes, espases, bastos}\}$\\[4pt]
Boles de l'1 al 20: \ «múltiple de 5» $=\{5,10,15,20\}$ \qquad «més gran que 20» $=\varnothing$
\end{exemple}

\vfill

\vfill
% ============================================================ 11.2
\clearpage
\seccio{sec2}{Esdeveniments}{Unió, intersecció i contrari.}

\vfill
\begin{tcbraster}[raster columns=4,raster equal height,raster column skip=6pt]
\begin{tcolorbox}[fmbase,colback=fmdefbg,colframe=fmdef,coltitle=white,title={Unió $A\cup B$},before upper=\centering]
\venn{\fill[fmblue!45] \ca \cb;}{\draw[fmnavy] \ca \cb; \node[font=\scriptsize] at (-0.95,0) {$A$}; \node[font=\scriptsize] at (0.95,0) {$B$};}\par
\small $A$ \textbf{o} $B$
\end{tcolorbox}
\begin{tcolorbox}[fmbase,colback=fmdefbg,colframe=fmdef,coltitle=white,title={Intersecció $A\cap B$},before upper=\centering]
\venn{\begin{scope}\clip \ca; \fill[fmblue!45] \cb; \end{scope}}{\draw[fmnavy] \ca \cb; \node[font=\scriptsize] at (-0.95,0) {$A$}; \node[font=\scriptsize] at (0.95,0) {$B$};}\par
\small $A$ \textbf{i} $B$ alhora
\end{tcolorbox}
\begin{tcolorbox}[fmbase,colback=fmdefbg,colframe=fmdef,coltitle=white,title={Contrari $\overline A$},before upper=\centering]
\venn{\fill[fmblue!45] (-1.9,-1.3) rectangle (1.9,1.3); \fill[white] (-0.55,0) circle (1);}{\draw[fmnavy] \ca; \node[font=\scriptsize] at (-0.55,0) {$A$};}\par
\small \textbf{no} passa $A$
\end{tcolorbox}
\begin{tcolorbox}[fmbase,colback=fmdefbg,colframe=fmdef,coltitle=white,title={Incompatibles},before upper=\centering]
\venn{}{\draw[fmnavy,fill=fmblue!20] (-0.95,0) circle (0.75) (0.95,0) circle (0.75); \node[font=\scriptsize] at (-0.95,0) {$A$}; \node[font=\scriptsize] at (0.95,0) {$B$};}\par
\small $A\cap B=\varnothing$
\end{tcolorbox}
\end{tcbraster}
{\comic\color{fmmuted}\centering Diagrames de Venn: el rectangle és l'espai mostral $E$. Un esdeveniment i el seu contrari sempre són incompatibles.\par}

\vfill
\begin{idea}[Dues relacions útils]
\textbf{De Morgan}: \ $\overline A\cap\overline B=\overline{A\cup B}$ \ («cap dels dos») \qquad $\overline A\cup\overline B=\overline{A\cap B}$ \ («no tots dos»)\\[3pt]
$A=(A\cap B)\cup(A\cap\overline B)$, \ dos trossos incompatibles.
\end{idea}

\vfill
\begin{exemple}[Exemple: \ dau, \ $A=$ «parell» $=\{2,4,6\}$, \ $B=$ «més de 3» $=\{4,5,6\}$]
$A\cup B=\{2,4,5,6\}$ \qquad $A\cap B=\{4,6\}$ \qquad $\overline A=\{1,3,5\}$ \qquad $\overline B=\{1,2,3\}$\\[4pt]
\textbf{No} són incompatibles: tenen en comú el 4 i el 6.
\end{exemple}

\vfill
\begin{exemple}[Exemple: \ boles de l'1 al 12]
«senar» $=\overline{\text{«parell»}}$ \qquad $D=\{6,12\}\subset C=\{3,6,9,12\}$: si passa $D$, passa $C$ \qquad $C\cap D=D$
\end{exemple}

\vfill
\begin{exemple}[Exemples]
Baralla espanyola, $O=$ «oros», $F=$ «figura»: \ $O\cap F$ = sota, cavall i rei d'oros (3) \qquad $O\cup F$: \ $10+12-3=19$ cartes\\[4pt]
«parell» i «un 5»: \textbf{incompatibles} \qquad «parell» i «més de 4»: \textbf{no} (tots dos tenen el 6)
\end{exemple}

\vfill

\vfill
% ============================================================ 11.3
\clearpage
\seccio{sec3}{Probabilitat d'un esdeveniment}{Un nombre entre 0 i 1 que diu com d'esperable és.}

\vfill
\begin{idea}[Freqüències]
Repetim un experiment $n$ vegades i $A$ ocorre $f$ vegades: $f$ és la \textbf{freqüència absoluta} i \ $h(A)=\dfrac{f}{n}$ \ la \textbf{freqüència relativa}.
\end{idea}

\vfill
\begin{center}\small
\taularod{>{\heavy}l<{\hspace{3pt}}*{6}{>{\hspace{3pt}}c<{\hspace{3pt}}}>{\hspace{3pt}}c<{\hspace{3pt}}}{
\rowcolor{fmblue}\color{white}Dau, 60 llançaments & \color{white}1 & \color{white}2 & \color{white}3 & \color{white}4 & \color{white}5 & \color{white}6 & \color{white}Total\\
Freq. absoluta $f$ & 8 & 11 & 9 & 12 & 10 & 10 & 60\\
\rowcolor{fmsurface}Freq. relativa $h$ & $0{,}13$ & $0{,}18$ & $0{,}15$ & $0{,}20$ & $0{,}17$ & $0{,}17$ & 1
}
\end{center}
{\centering\comic\color{fmmuted}freqüència relativa de «parell»: \ $\frac{11+12+10}{60}=0{,}55$\par}

\vfill
\begin{center}
\begin{tikzpicture}[x=1.45cm]
  \shade[left color=fmR!45,right color=fmK!55] (0,0) rectangle (10,0.45);
  \draw[fmnavy,line width=1pt] (0,0) rectangle (10,0.45);
  \foreach \x/\v/\t in {0/0/impossible,2.5/{0,25}/poc probable,5/{0,5}/tan probable com no,7.5/{0,75}/molt probable,10/1/segur} {
    \draw[fmnavy,line width=1pt] (\x,0) -- (\x,-0.15); \node[below,font=\small\heavy] at (\x,-0.15) {$\v$}; \node[below=14pt,font=\scriptsize,align=center,text width=2cm] at (\x,-0.15) {\t};}
\end{tikzpicture}
\end{center}

\vfill
\begin{minipage}[t]{0.52\linewidth}
\begin{idea}[Propietats]
\faCheck\ $0\le P(A)\le1$\\[2pt]
\faCheck\ $P(E)=1$ \ i \ $P(\varnothing)=0$\\[2pt]
\faCheck\ Les probabilitats dels resultats elementals \textbf{sumen 1}.\\[2pt]
\faCheck\ $P(\text{compost})$ = suma de les dels seus resultats:\\ \hspace*{1em}$P(\text{parell})=P(2)+P(4)+P(6)$
\end{idea}
\end{minipage}\hfill
\begin{minipage}[t]{0.45\linewidth}
\begin{truc}[Com la calculem?]
Si els resultats són \textbf{equiprobables}: llei de \textbf{Laplace} (11.4).\\[3pt]
Si no: amb la \textbf{freqüència relativa} de moltes repeticions (pàgina següent).
\end{truc}
\end{minipage}

\vfill
\begin{exemple}[Exemple: \ dau trucat amb $P(1)=0{,}1$, $P(2)=0{,}2$, $P(3)=0{,}1$, $P(4)=0{,}2$, $P(5)=0{,}1$]
$P(6)=1-(0{,}1+0{,}2+0{,}1+0{,}2+0{,}1)=\boxed{0{,}3}$ \qquad $P(\text{parell})=0{,}2+0{,}2+0{,}3=\boxed{0{,}7}$
\end{exemple}

\vfill
\begin{exemple}[Exemple: \ ruleta amb $P(\text{roig})=0{,}35$, $P(\text{blau})=0{,}25$, $P(\text{verd})=0{,}1$ i groc]
$P(\text{groc})=1-0{,}7=\boxed{0{,}3}$ \qquad $P(\text{blau o groc})=0{,}25+0{,}3=\boxed{0{,}55}$
\end{exemple}

\vfill

\vfill
\clearpage
\subsection{Llei dels grans nombres}

\vfill
\begin{idea}[Llei dels grans nombres]
Si repetim un experiment \textbf{moltes vegades}, la freqüència relativa s'estabilitza al voltant d'un valor: la \textbf{probabilitat}.\\[3pt]
{\centering\Large$P(A)\approx h(A)$ \ \normalsize si $n$ és gran\par}
\end{idea}

\vfill
\begin{minipage}[c]{0.5\linewidth}\centering\small
\taularod{>{\heavy}l<{\hspace{3pt}}*{6}{>{\hspace{3pt}}c<{\hspace{3pt}}}}{
\rowcolor{fmblue}\color{white}Xinxeta: $n$ & \color{white}10 & \color{white}50 & \color{white}100 & \color{white}200 & \color{white}500 & \color{white}1000\\
Punta amunt & 7 & 32 & 61 & 126 & 309 & 621\\
\rowcolor{fmsurface}$h$ & $0{,}7$ & $0{,}64$ & $0{,}61$ & $0{,}63$ & $0{,}618$ & $0{,}621$
}
\end{minipage}\hfill
\begin{minipage}[c]{0.46\linewidth}\centering
\begin{tikzpicture}[x=0.95cm,y=9cm]
  \draw[fmmuted!25] (0,0.5) grid[ystep=0.05,xstep=1] (6,0.75);
  \draw[-{Latex[length=2mm]},fmmuted,line width=0.8pt] (0,0.5) -- (6.4,0.5) node[right,font=\scriptsize] {$n$};
  \draw[-{Latex[length=2mm]},fmmuted,line width=0.8pt] (0,0.5) -- (0,0.78) node[above,font=\scriptsize] {$h$};
  \draw[fmR,dashed,line width=1pt] (0,0.62) -- (6,0.62) node[right,font=\scriptsize] {0,62};
  \draw[fmA,line width=1.6pt] (0.5,0.7) -- (1.5,0.64) -- (2.5,0.61) -- (3.5,0.63) -- (4.5,0.618) -- (5.5,0.621);
  \foreach \x/\y in {0.5/0.7,1.5/0.64,2.5/0.61,3.5/0.63,4.5/0.618,5.5/0.621} \fill[fmA] (\x,\y) circle (2.5pt);
  \foreach \x/\n in {0.5/10,2.5/100,5.5/1000} \node[below,font=\scriptsize] at (\x,0.5) {\n};
\end{tikzpicture}
\end{minipage}

\vfill
{\large Al principi la $h$ canvia molt; després s'acosta a $0{,}62$: \ $P(\text{punta amunt})\approx0{,}62$ \ i \ $P(\text{de costat})\approx0{,}38$.}

\vfill
\begin{tcbraster}[raster columns=2,raster equal height,raster column skip=10pt]
\begin{tcolorbox}[fmbase,colback=fmK!8,colframe=fmK,coltitle=white,title={\faBullseye\ \ Fer previsions}]
Si repetim $n$ vegades, esperem que $A$ ocórrega unes \ $\boxed{n\cdot P(A)}$ \ vegades.\par\smallskip
600 daus: \ $600\cdot\frac16=100$ sisos \textbf{aproximadament}.
\end{tcolorbox}
\begin{tcolorbox}[fmbase,colback=fmR!6,colframe=fmR,coltitle=white,title={\faBrain\ \ L'atzar no té memòria}]
Després de cinc creus seguides, la probabilitat de cara continua sent $\frac12$. La llei parla de moltíssimes repeticions, no de «compensar».
\end{tcolorbox}
\end{tcbraster}

\vfill
\begin{exemple}[Exemples]
\begin{minipage}[t]{0.48\linewidth}
Encerta 168 de 240 tirs lliures:\\
$P\approx\frac{168}{240}=0{,}7$ \ $\to$ \ en 50 tirs, uns $\boxed{35}$\\[5pt]
Xinxeta 500 vegades: de costat unes $500\cdot0{,}38=\boxed{190}$
\end{minipage}\hfill
\begin{minipage}[t]{0.49\linewidth}
20 boles; en 500 extraccions (tornant-la), 198 roges:\\
$\frac{x}{20}\approx\frac{198}{500}\approx0{,}4$ \ $\to$ \ unes $\boxed{8\text{ roges}}$\\[5pt]
Dau amb 502 sisos de 1000: \ esperaríem uns 167 \ $\to$ \textbf{trucat}
\end{minipage}
\end{exemple}

\vfill
\begin{truc}[Buffon i Pearson]
Buffon va llançar una moneda 4040 vegades: 2048 cares ($h\approx0{,}507$). Pearson, 24\,000 vegades: 12\,012 cares ($h=0{,}5005$). Totes dues s'acosten a $\frac12$.
\end{truc}

\vfill

\vfill
% ============================================================ 11.4
\clearpage
\seccio{sec4}{Càlcul de probabilitats}{Si tots els resultats són igual de probables, comptem.}

\vfill
\begin{minipage}[c]{0.6\linewidth}
\begin{tcolorbox}[enhanced,colback=fmyellow!18,colframe=fmgold,arc=4mm,boxrule=1.6pt,before upper=\centering,top=3mm,bottom=3mm]
{\Large$P(A)=\dfrac{\text{casos favorables}}{\text{casos possibles}}$}\par\smallskip
{\small només en experiments \textbf{regulars}: resultats equiprobables}
\end{tcolorbox}
\end{minipage}\hfill
\begin{minipage}[c]{0.37\linewidth}\centering
\foreach \i in {1,...,5}{\bola{fmR}\,}\\[2pt]
\foreach \i in {1,...,3}{\bola{fmA}\,}\ \bola{fmK}\,\bola{fmK}\\[3pt]
{\comic\small\color{fmmuted}5 roges, 3 blaves, 2 verdes}
\end{minipage}

\vfill
\begin{tcbraster}[raster columns=3,raster equal height,raster column skip=8pt]
\begin{tcolorbox}[fmbase,colback=fmdefbg,colframe=fmdef,coltitle=white,title={\faDiceSix\ \ Dau},before upper=\centering]
$P(\text{parell})=\dfrac36=\dfrac12$\par\smallskip
$P(\text{múltiple de 3})=\dfrac26=\dfrac13$
\end{tcolorbox}
\begin{tcolorbox}[fmbase,colback=fmdefbg,colframe=fmdef,coltitle=white,title={\faShoppingBag\ \ Bossa},before upper=\centering]
$P(R)=\dfrac{5}{10}=\dfrac12$\par\smallskip
$P(\text{no verda})=\dfrac{8}{10}=\dfrac45$
\end{tcolorbox}
\begin{tcolorbox}[fmbase,colback=fmdefbg,colframe=fmdef,coltitle=white,title={\faIcon{clone}\ \ Baralla (40)},before upper=\centering]
$P(\text{oros})=\dfrac{10}{40}=\dfrac14$\par\smallskip
$P(\text{figura})=\dfrac{12}{40}=\dfrac{3}{10}$
\end{tcolorbox}
\end{tcbraster}

\vfill
\begin{compte}[Només si són equiprobables]
Una ruleta amb sectors diferents, un dau trucat o una xinxeta \textbf{no} són regulars. «Aprovar o suspendre» té dos resultats, però no té per què ser $\frac12$ cadascun!
\end{compte}

\vfill
\begin{exemple}[Exemples]
\begin{minipage}[t]{0.48\linewidth}
Urna amb 4 roges, 6 blaves i 5 verdes (15):\\
$P(R)=\frac{4}{15}$, \ $P(B)=\frac{6}{15}=\frac25$, \ $P(\text{no }B)=\frac{9}{15}=\frac35$\\[5pt]
Dau: \ $P(\text{primer})=\frac36=\frac12$ \ \nota{2, 3, 5}
\end{minipage}\hfill
\begin{minipage}[t]{0.49\linewidth}
Lletres de PROBABILITAT (12 targetes):\\
$P(A)=\frac{2}{12}=\frac16$, \ $P(\text{vocal})=\frac{5}{12}$\\[5pt]
Copes \textbf{o} figura: \ $10+12-3=19$ \ \nota{sense repetir}\\
$P=\boxed{\tfrac{19}{40}}$
\end{minipage}
\end{exemple}

\vfill
\begin{exemple}[Exemple: \ 3 boles roges i unes quantes blaves, amb $P(R)=\frac15$. Quantes blaves hi ha?]
$\dfrac3n=\dfrac15\ \to\ n=15$ boles \ $\to$ \ $15-3=\boxed{12\text{ blaves}}$\\[5pt]
Quantes roges cal afegir perquè $P(R)=\frac12$? \ $\dfrac{3+x}{15+x}=\dfrac12\ \to\ 6+2x=15+x\ \to\ \boxed{x=9}$
\end{exemple}

\vfill

\vfill
\clearpage
\subsection{Propietats de la probabilitat}

\vfill
\begin{tcbraster}[raster columns=2,raster equal height,raster column skip=10pt,raster row skip=8pt]
\begin{tcolorbox}[fmbase,colback=fmdefbg,colframe=fmdef,coltitle=white,title={Entre 0 i 1}]
$0\le P(A)\le1$ \qquad $P(E)=1$, \ $P(\varnothing)=0$
\end{tcolorbox}
\begin{tcolorbox}[fmbase,colback=fmdefbg,colframe=fmdef,coltitle=white,title={Contrari}]
$P(\overline A)=1-P(A)$
\end{tcolorbox}
\end{tcbraster}

\vfill
\begin{tcolorbox}[enhanced,colback=fmyellow!18,colframe=fmgold,arc=4mm,boxrule=1.6pt,top=2mm,bottom=2mm]
\begin{minipage}[c]{0.65\linewidth}
{\Large$P(A\cup B)=P(A)+P(B)-P(A\cap B)$}\par\smallskip
\small Restem la intersecció perquè està comptada dues vegades.\\
Incompatibles: \ $P(A\cup B)=P(A)+P(B)$
\end{minipage}\hfill
\begin{minipage}[c]{0.32\linewidth}\centering
\venn{\begin{scope}\clip \ca; \fill[fmR!40] \cb; \end{scope}}{\draw[fmnavy] \ca \cb; \node[font=\scriptsize] at (-0.95,0) {$A$}; \node[font=\scriptsize] at (0.95,0) {$B$};}
\end{minipage}
\end{tcolorbox}

\vfill
\begin{idea}[També]
$P(A)=P(A\cap B)+P(A\cap\overline B)$ \qquad $P(\overline A\cap\overline B)=1-P(A\cup B)$
\end{idea}

\vfill
\begin{exemple}[Exemple: \ dau, \ $A=$ «parell», \ $B=$ «més de 3»]
$P(A\cup B)=\dfrac36+\dfrac36-\dfrac26=\dfrac46=\boxed{\dfrac23}$ \qquad \nota{comprovació: $\{2,4,5,6\}$, 4 de 6}
\end{exemple}

\vfill
\begin{exemple}[Exemple: \ oros o figura]
$P=\dfrac{10}{40}+\dfrac{12}{40}-\dfrac3{40}=\boxed{\dfrac{19}{40}}$ \qquad Copes o as: \ $\dfrac{10}{40}+\dfrac4{40}-\dfrac1{40}=\dfrac{13}{40}$
\end{exemple}

\vfill
\begin{exemple}[Exemple: \ $P(A)=0{,}5$, \ $P(\overline A\cap B)=0{,}3$, \ $P(A\cap B)=0{,}2$]
$P(B)=P(A\cap B)+P(\overline A\cap B)=0{,}2+0{,}3=\boxed{0{,}5}$ \qquad $P(A\cup B)=0{,}5+0{,}5-0{,}2=\boxed{0{,}8}$
\end{exemple}

\vfill
\begin{exemple}[Exemples]
$P(A)=0{,}6$, $P(B)=0{,}3$, $P(A\cap B)=0{,}1$: \ $P(A\cup B)=0{,}8$, \ $P(\overline A)=0{,}4$, \ $P(\overline A\cap\overline B)=0{,}2$\\[4pt]
$P(A\cup B)=0{,}7$, $P(A)=0{,}5$, $P(B)=0{,}4$: \ $P(A\cap B)=0{,}2$\\[4pt]
$P(A)=0{,}6$ i $P(B)=0{,}7$ \textbf{no} poden ser incompatibles: la unió valdria $1{,}3>1$
\end{exemple}

\vfill

\vfill
% ============================================================ 11.5
\clearpage
\seccio{sec5}{Experiències compostes. Probabilitat}{Diversos experiments seguits: arbre o taula.}

\vfill
\begin{idea}[Experiments compostos]
Estan formats per diversos experiments simples: dues monedes, una moneda i un dau, dos daus\dots\ \ El nombre de resultats és el \textbf{producte}: \ moneda i dau, $2\cdot6=12$; \ dos daus, $6\cdot6=36$.
\end{idea}

\vfill
\begin{minipage}[c]{0.45\linewidth}\centering\small
\providecommand{\sv}{}\renewcommand{\sv}{\cellcolor{fmyellow!70}\heavy 7}
\renewcommand{\arraystretch}{1.25}
\taularod{>{\columncolor{fmblue}\color{white}\heavy}c*{6}{w{c}{0.55cm}}}{
\rowcolor{fmblue}$+$ & \color{white}1 & \color{white}2 & \color{white}3 & \color{white}4 & \color{white}5 & \color{white}6\\
1 & 2 & 3 & 4 & 5 & 6 & \sv\\
2 & 3 & 4 & 5 & 6 & \sv & 8\\
3 & 4 & 5 & 6 & \sv & 8 & 9\\
4 & 5 & 6 & \sv & 8 & 9 & 10\\
5 & 6 & \sv & 8 & 9 & 10 & 11\\
6 & \sv & 8 & 9 & 10 & 11 & 12
}
\end{minipage}\hfill
\begin{minipage}[c]{0.52\linewidth}
\textbf{Dos daus}: 36 resultats equiprobables.\\[4pt]
$P(\text{suma }7)=\dfrac{6}{36}=\dfrac16$\qquad $P(\text{suma }2)=\dfrac{1}{36}$\\[6pt]
\begin{compte}[Les sumes no són equiprobables]
El 7 és la més probable. Amb $E=\{2,\dots,12\}$ i Laplace eixiria $\frac1{11}$: \mal{\text{incorrecte}}
\end{compte}
\end{minipage}

\vfill
\begin{minipage}[c]{0.5\linewidth}
\begin{idea}[Regla del producte]
Si els experiments són \textbf{independents}:\\[3pt]
\pas{1} En cada branca, la seua probabilitat.\\[2pt]
\pas{2} Un camí: \textbf{multipliquem} les branques.\\[2pt]
\pas{3} Diversos camins: \textbf{sumem}.
\end{idea}
\end{minipage}\hfill
\begin{minipage}[c]{0.47\linewidth}\centering
\begin{tikzpicture}[arbre,level 1/.style={sibling distance=1.7cm},level 2/.style={sibling distance=0.85cm},level distance=1.9cm]
  \node {}
    child {node {$F$} child {node {$F$} edge from parent node[ram,below] {0,3}} child {node {$E$} edge from parent node[ram,above] {0,7}} edge from parent node[ram,below] {0,3}}
    child {node {$E$} child {node {$F$} edge from parent node[ram,below] {0,3}} child {node {$E$} edge from parent node[ram,above] {0,7}} edge from parent node[ram,above] {0,7}};
  \node[font=\scriptsize,text=fmK!80!black] at (4.6,1.27) {$0{,}49$}; \node[font=\scriptsize] at (4.6,0.42) {$0{,}21$};
  \node[font=\scriptsize] at (4.6,-0.42) {$0{,}21$}; \node[font=\scriptsize,text=fmR] at (4.6,-1.27) {$0{,}09$};
\end{tikzpicture}\\
{\comic\small\color{fmmuted}dos tirs lliures al 70\,\%}
\end{minipage}

\vfill
\begin{exemple}[Exemple: \ dos tirs lliures, encerta el 70\,\% ($E$ encert, $F$ fallada)]
Tots dos: \ $0{,}7\cdot0{,}7=\boxed{0{,}49}$ \qquad Cap: \ $0{,}3\cdot0{,}3=\boxed{0{,}09}$ \qquad Exactament un: \ $0{,}21+0{,}21=\boxed{0{,}42}$
\end{exemple}

\vfill
\begin{exemple}[Exemples]
\begin{minipage}[t]{0.48\linewidth}
Moneda i dau: \ $P(C\text{ i }6)=\frac12\cdot\frac16=\boxed{\tfrac{1}{12}}$\\[4pt]
Tres monedes, exactament dues cares:\\ 3 camins de $\frac18$: \ $\boxed{\tfrac38}$; \ almenys dues: $\boxed{\tfrac48=\tfrac12}$
\end{minipage}\hfill
\begin{minipage}[t]{0.49\linewidth}
Dos daus, suma 10: \ $\frac{3}{36}=\boxed{\tfrac1{12}}$\\[3pt]
Suma 10 o més: \ $\frac{3+2+1}{36}=\boxed{\tfrac16}$\\[3pt]
Doble: \ $\frac{6}{36}=\boxed{\tfrac16}$ \quad Almenys un 6: \ $\boxed{\tfrac{11}{36}}$
\end{minipage}
\end{exemple}

\vfill

\vfill
% ============================================================ 11.6
\clearpage
\seccio{sec6}{Experiències dependents. Probabilitat condicionada}{Tornem la bola a la bossa, o no?}

\vfill
{\large Traiem dues boles d'una bossa amb \textbf{3 roges} \bola{fmR}\bola{fmR}\bola{fmR} i \textbf{2 blaves} \bola{fmA}\bola{fmA}.}

\vfill
\begin{tcbraster}[raster columns=2,raster equal height,raster column skip=10pt]
\begin{tcolorbox}[fmbase,colback=fmK!8,colframe=fmK,coltitle=white,title={Amb reemplaçament: independents},before upper=\centering]
\begin{tikzpicture}[arbre,level 1/.style={sibling distance=1.5cm},level 2/.style={sibling distance=0.75cm},level distance=1.7cm]
  \node {}
    child {node {$B$} child {node {$B$} edge from parent node[ram,below] {$\frac25$}} child {node {$R$} edge from parent node[ram,above] {$\frac35$}} edge from parent node[ram,below] {$\frac25$}}
    child {node {$R$} child {node {$B$} edge from parent node[ram,below] {$\frac25$}} child {node {$R$} edge from parent node[ram,above] {$\frac35$}} edge from parent node[ram,above] {$\frac35$}};
\end{tikzpicture}\par\smallskip
$P(RR)=\dfrac35\cdot\dfrac35=\boxed{\dfrac{9}{25}}$
\end{tcolorbox}
\begin{tcolorbox}[fmbase,colback=fmR!6,colframe=fmR,coltitle=white,title={Sense reemplaçament: dependents},before upper=\centering]
\begin{tikzpicture}[arbre,level 1/.style={sibling distance=1.5cm},level 2/.style={sibling distance=0.75cm},level distance=1.7cm]
  \node {}
    child {node {$B$} child {node {$B$} edge from parent node[ram,below] {$\frac14$}} child {node {$R$} edge from parent node[ram,above] {$\frac34$}} edge from parent node[ram,below] {$\frac25$}}
    child {node {$R$} child {node {$B$} edge from parent node[ram,below] {$\frac24$}} child {node {$R$} edge from parent node[ram,above] {$\frac24$}} edge from parent node[ram,above] {$\frac35$}};
\end{tikzpicture}\par\smallskip
$P(RR)=\dfrac35\cdot\dfrac24=\boxed{\dfrac{3}{10}}$
\end{tcolorbox}
\end{tcbraster}

\vfill
\begin{truc}[Sense reemplaçament]
En la segona branca pregunta't: «quantes d'aquest color queden?» i «quantes queden en total?». Si la primera ha sigut roja, queden 2 roges de 4 boles.
\end{truc}

\vfill
\begin{minipage}[c]{0.42\linewidth}\raggedright
\textbf{Com resoldre un problema}:\\[4pt]
\pas{1} Quin és l'experiment? Simple o compost? Equiprobable?\\[3pt]
\pas{2} Simple i regular: \textbf{Laplace}. Una taula: \textbf{comptem}.\\[3pt]
\pas{3} Compost: \textbf{arbre}, multipliquem i sumem camins.
\end{minipage}\hfill
\begin{minipage}[c]{0.55\linewidth}
\begin{exemple}[Exemple: \ dues cartes sense tornar-les. Dos oros?]
$P=\dfrac{10}{40}\cdot\dfrac{9}{39}=\dfrac{90}{1560}=\boxed{\dfrac{3}{52}}\approx0{,}06$
\end{exemple}
\begin{exemple}[Exemple: \ moneda i dues urnes]
Urna I (2 R, 3 B) si cara; urna II (4 R, 1 B) si creu:\\[2pt]
$P(R)=\frac12\cdot\frac25+\frac12\cdot\frac45=\frac15+\frac25=\boxed{\tfrac35}$
\end{exemple}
\end{minipage}

\vfill
\begin{exemple}[Exemple: \ en una classe de 30, triem una persona a l'atzar]
\begin{minipage}[c]{0.42\linewidth}\centering\small
\taularod[fmex]{>{\heavy}l<{\hspace{4pt}}>{\hspace{4pt}}c<{\hspace{4pt}}>{\hspace{4pt}}c<{\hspace{4pt}}>{\hspace{4pt}}c<{\hspace{4pt}}}{
\rowcolor{fmex} & \color{white}\heavy Ulleres & \color{white}\heavy Sense & \color{white}\heavy Total\\
Xiques & 6 & 10 & 16\\
\rowcolor{fmexbg}Xics & 4 & 10 & 14\\
Total & 10 & 20 & 30
}
\end{minipage}\hfill
\begin{minipage}[c]{0.55\linewidth}
$P(\text{ulleres})=\frac{10}{30}=\frac13$ \qquad $P(\text{xica i ulleres})=\frac{6}{30}=\frac15$\\[3pt]
$P(\text{xica o ulleres})=\frac{16+10-6}{30}=\frac23$\\[3pt]
Si és un xic: \ $P(\text{ulleres})=\frac{4}{14}=\frac27$ \ \nota{només 14 casos}
\end{minipage}
\end{exemple}

\vfill

\vfill
\clearpage
\subsection{Problemes amb arbres de dues etapes}

\vfill
\begin{tcbraster}[raster columns=2,raster equal height,raster column skip=10pt]
\begin{tcolorbox}[fmbase,colback=fmK!8,colframe=fmK,coltitle=white,title={Probabilitat total}]
Sumem \textbf{tots els camins} que porten a $B$:\\[3pt]
\small$P(B)=P(A_1)\,P(B\cond A_1)+P(A_2)\,P(B\cond A_2)$
\end{tcolorbox}
\begin{tcolorbox}[fmbase,colback=fmA!8,colframe=fmA,coltitle=white,title={A posteriori (Bayes)}]
Si ha eixit $B$, la probabilitat que vinga d'$A_1$:\\[3pt]
$P(A_1\cond B)=\dfrac{\text{camí d'}A_1\text{ a }B}{\text{tots els camins a }B}$
\end{tcolorbox}
\end{tcbraster}

\vfill
\begin{exemple}[Exemple: \ una prova mèdica]
{\small 2\,\% de malalts. La prova dona positiu en el 95\,\% dels malalts i en el 4\,\% dels sans.}\par\smallskip
\begin{minipage}[c]{0.5\linewidth}\centering
\begin{tikzpicture}[x=1.3cm,y=0.5cm]
  \node[nd] (r) at (0,0) {};
  \node[nd,fill=fmR,minimum size=1.6em] (m) at (1.6,2) {M}; \node[nd,fill=fmK,minimum size=1.6em] (s) at (1.6,-2) {S};
  \node[nd,fill=fmsf] (mp) at (3.3,3) {$+$}; \node[nd,fill=fmmuted] (mn) at (3.3,1) {$-$};
  \node[nd,fill=fmsf] (sp) at (3.3,-1) {$+$}; \node[nd,fill=fmmuted] (sn) at (3.3,-3) {$-$};
  \draw[br] (r)--node[pl]{$0{,}02$}(m); \draw[br] (r)--node[pl]{$0{,}98$}(s);
  \draw[br] (m)--node[pl]{$0{,}95$}(mp); \draw[br] (m)--node[pl]{$0{,}05$}(mn);
  \draw[br] (s)--node[pl]{$0{,}04$}(sp); \draw[br] (s)--node[pl]{$0{,}96$}(sn);
  \node[right=3pt,font=\small\heavy,text=fmsf] at (mp) {$0{,}019$}; \node[right=3pt,font=\small\heavy,text=fmsf] at (sp) {$0{,}0392$};
\end{tikzpicture}
\end{minipage}\hfill
\begin{minipage}[c]{0.48\linewidth}
$P(M\cap+)=0{,}02\cdot0{,}95=0{,}019$\\[4pt]
$P(+)=0{,}019+0{,}0392=\boxed{0{,}0582}$\\[4pt]
$P(M\cond+)=\dfrac{0{,}019}{0{,}0582}\approx\boxed{0{,}33}$
\end{minipage}
\end{exemple}

\vfill
\begin{compte}[Sorprenent!]
Només 1 de cada 3 positius està malalt de veritat: la malaltia és rara i hi ha molts més sans que donen positiu per error.
\end{compte}

\vfill
\begin{exemple}[Exemple: \ urna I (2 R, 3 B, 5 N) i urna II (4 R, 2 B, 2 N); triem urna a l'atzar]
$P(\text{blanca})=\dfrac12\cdot\dfrac3{10}+\dfrac12\cdot\dfrac28=\dfrac3{20}+\dfrac18=\boxed{\dfrac{11}{40}}$ \qquad $P(\text{I}\cond\text{blanca})=\dfrac{3/20}{11/40}=\boxed{\dfrac6{11}}$
\end{exemple}

\vfill
\begin{exemple}[Exemples]
Moneda: cara, bossa A (2 R, 3 B); creu, bossa B (4 R, 1 B): \ \\ \hspace*{1em}$P(R)=\frac12\cdot\frac25+\frac12\cdot\frac45=\frac35$, \quad $P(A\cond R)=\frac{1/5}{3/5}=\frac13$\\[4pt]
70\,\% aprova Mates; d'ells el 80\,\% aprova Física, i dels altres el 30\,\%:\\ \hspace*{1em}$P(F)=0{,}7\cdot0{,}8+0{,}3\cdot0{,}3=0{,}65$, \quad $P(M\cond F)=\frac{0{,}56}{0{,}65}\approx0{,}86$
\end{exemple}

\vfill

\vfill
% ============================================================ 11.7
\clearpage
\seccio{sec7}{Taules de contingència}{La condició és el nou espai mostral.}

\vfill
\begin{minipage}[c]{0.52\linewidth}\centering\small
{\heavy 100 alumnes de 4t: practiquen esport?}\par\smallskip
\taularod{>{\heavy}l*{3}{>{\hspace{5pt}}c<{\hspace{5pt}}}}{%
\rowcolor{fmblue} & \color{white}Esport $S$ & \color{white}No $\overline S$ & \color{white}Total\\
Xics $H$ & $28$ & $12$ & $40$\\
\rowcolor{fmyellow!30}Xiques $D$ & $33$ & $27$ & $60$\\
Total & $61$ & $39$ & $100$
}
\end{minipage}\hfill
\begin{minipage}[c]{0.45\linewidth}
\begin{idea}[Taula de contingència]
Classifica segons dues característiques. Triant a l'atzar, Laplace amb els totals:\\[2pt]
$P(S)=0{,}61$, \ $P(D)=0{,}6$, \ $P(D\cap S)=0{,}33$
\end{idea}
\end{minipage}

\vfill
\begin{tcolorbox}[enhanced,colback=fmyellow!18,colframe=fmgold,arc=4mm,boxrule=1.6pt,top=2mm,bottom=2mm]
\begin{minipage}[c]{0.4\linewidth}\centering
{\Large$P(A\cond B)=\dfrac{P(A\cap B)}{P(B)}$}\par\smallskip
\small «$A$ sabent que ha passat $B$»
\end{minipage}\hfill
\begin{minipage}[c]{0.57\linewidth}
En una taula: ens quedem amb la fila o columna de $B$ i hi apliquem Laplace.\\[3pt]
\textbf{Independents}: \ $P(A\cond B)=P(A)$, \ és a dir, $P(A\cap B)=P(A)\cdot P(B)$.
\end{minipage}
\end{tcolorbox}

\vfill
\begin{exemple}[Exemple: \ amb la taula dels 100 alumnes]
Una xica fa esport: fila de les xiques, \ $P(S\cond D)=\dfrac{33}{60}=\boxed{0{,}55}$\\[4pt] Un esportista és xica: \ $P(D\cond S)=\dfrac{33}{61}\approx0{,}54$\\[4pt]
$P(S\cond D)=0{,}55\neq P(S)=0{,}61$: \ \textbf{dependents} \nota{(o bé $0{,}6\cdot0{,}61=0{,}366\neq0{,}33$)}
\end{exemple}

\vfill
\begin{exemple}[Exemple: \ 150 viatgers: 60 parlen anglés $I$, 45 francés $F$, 15 tots dos]
\begin{minipage}[c]{0.42\linewidth}\centering\small
\taularod{>{\heavy}l*{3}{>{\hspace{4pt}}c<{\hspace{4pt}}}}{%
\rowcolor{fmblue} & \color{white}$F$ & \color{white}$\overline F$ & \color{white}Total\\
$I$ & $15$ & $45$ & $60$\\
\rowcolor{fmsurface}$\overline I$ & $30$ & $60$ & $90$\\
Total & $45$ & $105$ & $150$
}
\end{minipage}\hfill
\begin{minipage}[c]{0.56\linewidth}
$P(I\cup F)=\dfrac{60+45-15}{150}=0{,}6$\\[4pt]
$P(F\cond I)=\dfrac{15}{60}=0{,}25$\\[4pt]
Només francés: \ $P(F\cap\overline I)=\dfrac{30}{150}=0{,}2$
\end{minipage}
\end{exemple}

\vfill
\begin{exemple}[Exemple: \ sense taula, \ $P(A)=0{,}5$, \ $P(B)=0{,}4$, \ $P(A\cup B)=0{,}7$]
$P(A\cap B)=0{,}5+0{,}4-0{,}7=0{,}2$ \qquad $P(A\cond B)=\dfrac{0{,}2}{0{,}4}=0{,}5=P(A)$: \ \textbf{independents}
\end{exemple}

\vfill

\vfill
% ============================================================ resum
\resumtitol
\begin{tcbraster}[raster columns=2,raster equal height=rows,raster column skip=14pt,raster row skip=34pt,raster before skip=18pt]
\rtile{11.1}{Experiments}{Aleatori: no sabem el resultat.\\[1ex] Espai mostral $E$: tots els resultats.}
\rtile{11.2}{Esdeveniments}{$A\cup B$, $A\cap B$, $\overline A$. Incompatibles: $A\cap B=\varnothing$.}
\rtile{11.3}{Probabilitat}{Moltes repeticions: $P(A)\approx h(A)$.\par\vspace{1ex}{\centering$0\le P(A)\le1$\par}}
\rtile{11.4}{Laplace}{$P(A)=\frac{\text{favorables}}{\text{possibles}}$, si són equiprobables. $P(\overline A)=1-P(A)$.}
\rtile{11.5}{Compostes}{Arbre o taula.\\[1ex] Un camí: multipliquem.\\[1ex] Diversos camins: sumem.}
\rtile{11.6}{Dependents}{Sense reemplaçament canvien les branques. Probabilitat total: sumem camins.}
\rtile{11.7}{Contingència}{$P(A\cond B)=\frac{P(A\cap B)}{P(B)}$: la fila o columna de $B$.}
\end{tcbraster}
\end{document}
